% =============================================================================
%  text/introduction.tex
%  Introduction: demonstrates citations, Chinese paragraph, and a bulleted
%  list of contributions. Replace all content with your own.
%  引言示例：演示引用、中英混排段落与要点列表。
% =============================================================================

A growing literature examines how political elites shape state capacity and
political survival \citep{author2018book,author2020example}. Scholars have
shown that cohesive elite networks can both strengthen the state and threaten
the ruler, a tension sometimes called the ruler's dilemma
\citep[e.g.,][]{author2019chapter}. This paper develops a dynamic framework
that explains how rulers' strategic choices shift across different political
conditions.

% 中文段落示例：中英文段落交错排布（英文为主、中文为辅）。
% 中文段落示例：中英文段落交错排布（英文为主、中文为辅）。
当统治者的权力稳固时，培养具有凝聚力的精英网络有助于提升国家能力；
当权力脆弱时，统治者则可能转而分化精英联盟以阻止挑战。
本文的理论框架将这两种策略统一在一个基于边际收益的决策模型之中，
并说明理性的统治者总是选择边际收益最大的策略。

This paper makes three contributions. First, it provides a unified framework
that connects elite management with state-building outcomes. Second, it
introduces a novel empirical strategy for measuring elite network structure.
Third, it applies the framework to a salient historical case with rich
temporal variation in the ruler's power security.

Our empirical analysis proceeds as follows. Section~\ref{sec:elite-ruler}
presents the theoretical argument. Section~\ref{sec:case-background}
introduces the historical background of the case. Section~\ref{sec:sample}
describes the research design, and Section~\ref{sec:main-results} reports the
results. Section~\ref{sec:conclusion} concludes.

% 贡献列表示例（如需编号列表，请将 itemize 改为 enumerate）
Our main findings can be summarized as follows:
\begin{itemize}
  \item Well-connected elites enjoy better promotion prospects when the
        ruler's authority is consolidated.
  \item The same elites become targets of suppression when the ruler's
        authority is challenged.
  \item These dynamics have lasting consequences for state capacity.
\end{itemize}
